\documentclass[11pt]{article}
\usepackage[a4paper,margin=1in]{geometry}
\usepackage{amsmath,amssymb,amsthm}
\usepackage{graphicx}
\usepackage{hyperref}

\newtheorem{theorem}{Theorem}
\newtheorem{lemma}{Lemma}

\newcommand{\Om}{\Omega}
\newcommand{\R}{\mathbb R}
\newcommand{\ip}[2]{\langle #1,#2\rangle}
\newcommand{\fp}{\mathfrak p}

\title{A partial answer to the expected \(p=2\) rate in Chen--Yousept}
\author{FA Banach run \texttt{fa\_banach\_001}, lane 12}
\date{June 20, 2026}

\begin{document}
\maketitle

\section*{Status}

This packet records a partial positive answer to Remark 4.9 of Chen--Yousept
\cite{ChenYousept2020}. The source paper obtains, in the weakest-data case of
Theorem 4.6(iv), a rate which is suboptimal when the penalty exponent is
\(p=2\), and writes that the expected rate is
\[
  \|a^\delta_{\alpha(\delta)}-a^\dagger\|_{2}^{2}
  =O\!\left(\delta^{\frac{2s}{s+2}}\right).
\]
The result below proves this expected rate under a stronger true-state
multiplier assumption and an adapted Hilbert-scale smoothness assumption. The
varying coefficients remain in the source paper's original
\(L^{\mathfrak p}\)-bounded admissible class, with \(\mathfrak p\geq2\).

\begin{figure}[h]
  \centering
  \includegraphics[width=0.95\linewidth]{figures/open_problem_crop.png}
  \caption{Remark 4.9 in \cite{ChenYousept2020}, page 23.}
\end{figure}

\section*{Extracted question}

In the notation of \cite[Theorem 4.6(iv)]{ChenYousept2020}, for the Hilbert
penalty \(p=2\), can one recover the expected rate
\[
  \|a^\delta_{\alpha(\delta)}-a^\dagger\|_2^2
  =O\!\left(\delta^{\frac{2s}{s+2}}\right)?
\]

\section*{Idea of the proof}

The loss in the source paper comes from estimating the coefficient difference
in \(H_q^1(\Omega)^*\) and then interpolating an \(H^2\)-bound for the state
difference with the weak observation norm. This introduces a square root in the
stability estimate before the variational-source-condition argument starts.

For \(p=2\) one can instead use the equation itself. Let
\[
  w:=S(a)-S(a^\dagger),\qquad h:=a^\dagger-a .
\]
If \(A_0\) is the fixed Neumann realization of
\(-\operatorname{div}(\kappa\nabla\cdot)+1\), then testing the difference
equation with \(\varphi\in D(A_0)\) gives
\[
  \int_\Omega h\,u^\dagger\varphi\,dx
  =\ip{w}{A_0\varphi}_{L^2}
    +\int_\Omega (a-1)w\varphi\,dx .
\]
Since \(D(A_0)\subset H^2(\Omega)\subset L^\infty(\Omega)\) in dimensions two
and three, the source paper's \(L^{\mathfrak p}\)-bound on \(a\), with
\(\mathfrak p\ge2\), is enough to control the nonlinear remainder by the
\(L^2\)-size of \(w\). Thus \(h\) is controlled linearly in the dual of the
adapted second-order space \(u^\dagger D(A_0)\). A standard Hilbert-scale
frequency splitting then gives exactly the exponent \(2s/(s+2)\).

\section*{Setting}

Let \(n\in\{2,3\}\). Let \(\Omega\subset\R^n\) be a bounded \(C^{1,1}\) domain
and let \(\kappa\in C^{0,1}(\overline\Omega)\) satisfy
\[
  0<\lambda_0\leq \kappa(x)\leq \Lambda<\infty .
\]
Assume that the admissible coefficients are those of the source paper:
\[
  \mathcal D
  :=\{a\in L^\fp(\Omega): \|a\|_{L^\fp}\leq M,\ 
       0<\underline a\leq a(x)\text{ a.e.}\},
  \qquad \fp\geq2 .
\]
For \(a\in\mathcal D\), let \(S(a)\) be the weak solution of the Neumann
problem considered in \cite[Section 4]{ChenYousept2020}, with fixed measure
data. Fix a true coefficient \(a^\dagger\in\mathcal D\) and write
\(u^\dagger:=S(a^\dagger)\). Assume
\[
  0<\gamma\leq |u^\dagger|\leq \Gamma<\infty
  \quad\text{a.e. on }\Omega .
\]

Let \(H=L^2(\Omega)\), and let \(A_0\) be the positive self-adjoint Neumann
realization in \(H\) of
\[
  -\operatorname{div}(\kappa\nabla\cdot)+1 .
\]
Thus
\[
  \ip{A_0\varphi,\psi}_H
  =\int_\Omega \kappa\nabla\varphi\cdot\nabla\psi+\varphi\psi\,dx,
  \qquad \varphi\in D(A_0),\ \psi\in H^1(\Omega).
\]
We use the standard \(H^2\)-regularity \(D(A_0)\hookrightarrow H^2(\Omega)\);
for \(n\le3\), this gives \(D(A_0)\hookrightarrow L^\infty(\Omega)\).
Define the adapted second-order space
\[
  X_2:=u^\dagger D(A_0),\qquad
  \|g\|_{X_2}:=\|A_0(g/u^\dagger)\|_2 .
\]
Since multiplication by \(u^\dagger\) is boundedly invertible on \(H\), \(X_2\)
is densely and continuously embedded in \(H\). Let \((X_t)_{-2\leq t\leq2}\)
be the Hilbert scale generated by \(X_2\hookrightarrow H\), and write
\(X_{-2}\) for the dual of \(X_2\) with pivot \(H\).

\begin{lemma}[Linear adapted \(X_{-2}\)-stability]\label{lem:stability}
There is a constant \(C\), depending only on the fixed ellipticity and
admissible-class bounds, such that for every \(a\in\mathcal D\),
\[
  \|a^\dagger-a\|_{X_{-2}}
  \leq C\,\|S(a)-S(a^\dagger)\|_2 .
\]
\end{lemma}

\begin{proof}
Put \(h=a^\dagger-a\) and \(w=S(a)-S(a^\dagger)\). For \(g\in X_2\), write
\(g=u^\dagger\varphi\) with \(\varphi\in D(A_0)\). Subtracting the weak
equations for \(S(a)\) and \(S(a^\dagger)\), tested against \(\varphi\), gives
\[
  \int_\Omega h\,u^\dagger\varphi\,dx
  =\int_\Omega \kappa\nabla w\cdot\nabla\varphi+a\,w\varphi\,dx .
\]
Adding and subtracting \(\int_\Omega w\varphi\,dx\) on the right side yields
\[
  \int_\Omega h\,g\,dx
  =\ip{w}{A_0\varphi}_H+\int_\Omega (a-1)w\varphi\,dx .
\]
Since \(\fp\ge2\), the conjugate exponent \(\fp'\le2\). By boundedness of
\(\Omega\), \(\|w\|_{\fp'}\leq C\|w\|_2\). Since \(\|a\|_{L^\fp}\le M\) and
\(D(A_0)\hookrightarrow L^\infty(\Omega)\), we get
\[
  \left|\int_\Omega (a-1)w\varphi\,dx\right|
  \leq C(1+M)\|w\|_2\|\varphi\|_\infty
  \leq C(1+M)\|w\|_2\|A_0\varphi\|_2 .
\]
The first term satisfies
\[
  |\ip{w}{A_0\varphi}_H|\leq \|w\|_2\|A_0\varphi\|_2 .
\]
Combining the two estimates and using
\(\|g\|_{X_2}=\|A_0\varphi\|_2\) proves
\[
  \left|\int_\Omega h\,g\,dx\right|
  \leq C\|w\|_2\|g\|_{X_2}.
\]
Taking the supremum over \(\|g\|_{X_2}\leq1\) proves the claim.
\end{proof}

\begin{lemma}[Hilbert-scale variational inequality]\label{lem:vsc}
Let \(0<s\leq2\) and \(f^\dagger:=a^\dagger-a^*\in X_s\). For every
\(\beta>0\) there is a constant \(C_\beta>0\) such that, for all
\(a\in\mathcal D\),
\[
  \ip{a^\dagger-a}{f^\dagger}_H
  \leq \frac{\beta}{2}\|a-a^\dagger\|_2^2
      +C_\beta\,
       \|S(a)-S(a^\dagger)\|_2^{\frac{2s}{s+2}} .
\]
\end{lemma}

\begin{proof}
Let \(G\geq I\) be the positive self-adjoint scale generator, so that
\(X_t=D(G^{t/2})\) for \(t\geq0\) and
\(\|z\|_{X_{-2}}=\|G^{-1}z\|_H\). Let \(E_\lambda\) be the spectral resolution
of \(G\), and set \(P_R=E_G([1,R])\) for \(R\geq1\). For
\(h:=a^\dagger-a\),
\[
  \ip{h}{f^\dagger}
  =\ip{h}{(I-P_R)f^\dagger}+\ip{h}{P_R f^\dagger}.
\]
The high-frequency term satisfies, by Young's inequality,
\[
  |\ip{h}{(I-P_R)f^\dagger}|
  \leq \|h\|_2 R^{-s/2}\|f^\dagger\|_{X_s}
  \leq \frac{\beta}{2}\|h\|_2^2
      +C_\beta R^{-s}\|f^\dagger\|_{X_s}^2 .
\]
The low-frequency term is bounded by
\[
  |\ip{h}{P_R f^\dagger}|
  \leq \|h\|_{X_{-2}}\|P_R f^\dagger\|_{X_2}
  \leq R^{(2-s)/2}\|h\|_{X_{-2}}\|f^\dagger\|_{X_s}.
\]
Using Lemma \ref{lem:stability}, the total upper bound is
\[
  \frac{\beta}{2}\|h\|_2^2
  +C_\beta R^{-s}
  +C R^{(2-s)/2}\|S(a)-S(a^\dagger)\|_2 .
\]
Optimizing in \(R\) gives
\[
  R\simeq \|S(a)-S(a^\dagger)\|_2^{-2/(s+2)}
\]
when the data discrepancy is small; for large discrepancies the constant can
be enlarged. Substitution yields the exponent \(2s/(s+2)\).
\end{proof}

\begin{theorem}[Expected \(p=2\) rate in the adapted scale]
Assume the setting above, let \(0<s\leq2\), and suppose
\(a^\dagger-a^*\in X_s\). Then the \(L^2\)-penalty Tikhonov regularization
problem satisfies a variational source condition with index
\[
  \Psi(t)=C t^{\frac{2s}{s+2}} .
\]
Consequently, for the parameter choice used in
\cite[Corollary 3.2]{ChenYousept2020},
\[
  \|a^\delta_{\alpha(\delta)}-a^\dagger\|_2^2
  =O\!\left(\delta^{\frac{2s}{s+2}}\right).
\]
The same conclusion holds for data norms \(Y=L^\tau(\Omega)\), \(\tau\geq2\),
because \(\Omega\) is bounded and \(\|v\|_2\leq C\|v\|_\tau\).
\end{theorem}

\begin{proof}
The variational source condition is Lemma \ref{lem:vsc}. The convergence-rate
claim is the standard implication from a variational source condition to
Tikhonov rates, applied exactly as in \cite[Corollary 3.2]{ChenYousept2020}.
\end{proof}

\section*{Verification and limitations}

The proof is complete for the stated strengthened setting. Compared with the
previous partial version, Lemma \ref{lem:stability} no longer uses a uniform
\(L^\infty\)-bound on the varying admissible coefficients; the source paper's
\(L^\fp\)-bound with \(\fp\ge2\) is enough. The result still does not settle
the full source-paper formulation, because the smoothness assumption is in the
adapted scale \(X_s\) generated by \(u^\dagger D(A_0)\), rather than directly
in the classical Triebel--Lizorkin/Sobolev scale of the source paper. In smooth
positive situations, this adapted scale is the natural Neumann elliptic Hilbert
scale and is equivalent to the corresponding Sobolev scale after multiplication
by \(u^\dagger\).

Bounded literature checks were performed for the exact source title, arXiv id,
the phrase ``\(2s/(s+2)\)'', and combinations of ``unbounded diffusion'',
``measure data'', ``variational source condition'', and ``Tikhonov''. No later
paper explicitly resolving Remark 4.9 was found in those searches. The
Hilbert-scale splitting itself is standard; the new point recorded here is the
short adapted \(X_{-2}\) stability estimate for this nonlinear coefficient map,
now with only the source paper's \(L^\fp\)-control of the varying coefficient.

\section*{Human-review recommendation}

First check Lemma \ref{lem:stability}; it is the only problem-specific step. If
accepted, the remaining Hilbert-scale argument is standard. A useful next step
would be to translate \(X_s\) into the exact Littlewood--Paley notation of
\cite{ChenYousept2020} under the weakest true-state multiplier assumptions.

\begin{thebibliography}{9}
\bibitem{ChenYousept2020}
D.-H. Chen and I. Yousept,
\emph{Variational source conditions in \(L^p\)-spaces},
arXiv:2007.13474, 2020.

\bibitem{HofmannMathe2012}
B. Hofmann and P. Math\'e,
\emph{Parameter choice in Banach space regularization under variational
inequalities},
Inverse Problems 28 (2012), 104006.
\end{thebibliography}

\end{document}
